\chapter{删了一张图，后面几十个编号全要手动改}

\begin{notebox}
本章解决 Word 的第三笔血债：交叉引用。学完这章，“如图 X 所示”永远跟着图走，删图加图你都不用回头管。
\end{notebox}

\painhead{烦恼}

“删了图 3，正文里十几处'如图 3 所示'全要找出来改成图 2，后面还有几十个编号顺延……改到怀疑人生。”

\painhead{现象 · 这些场景你肯定见过}

\textbf{场景一}：审稿意见要求删掉一张意义不大的图。你删了，然后开始全文搜索“图 4”改成“图 3”，“图 5”改成“图 4”……改完之后发现还有一处写着“如图 6(2)(b) 所示”，它的图早就不存在了。

\textbf{场景二}：参考文献编号灾难。你引用了 20 篇文献，正文里全是手打的 [12]。导师说再补两篇最新的——新文献按年份要插到第 5 条后面，于是 [6] 变 [8]，[7] 变 [9]……你手动顺延了两个小时，最后有一处 [13] 和 [14] 写反了，直接被退回。

\textbf{场景三}：公式编号。“式（3-5）”手打了十几处，在式（3-2）后面插入一个新公式，后半个章节的编号全要挪。

病根：\textbf{你在 Word 里手打的编号，是死的}。它是你抄上去的一个数字，图变了它不会跟着变。

LaTeX 的编号是\textbf{活的}：每张图、每个公式、每节标题都有编号，你给它们起代号，正文里用代号指路。图删了加了几号，代号不变，数字自动重算。

\painhead{解决线路}

\subsection*{5.1 核心二件套：\textbackslash{}label + \textbackslash{}ref}

上一章你已经见过 \texttt{\textbackslash{}label\{fig:result\}} 了，当时是配角，现在是主角。完整跑一遍：

\begin{lstlisting}
\documentclass[UTF8, 12pt]{ctexart}
\usepackage{graphicx}
\begin{document}

实验结果如图~\ref{fig:result} 所示，效果拔群。

\begin{figure}[htbp]
  \centering
  \includegraphics[width=0.8\textwidth]{result.png}
  \caption{实验结果对比图}
  \label{fig:result}
\end{figure}

数据解读详见第~\ref{sec:discuss} 节。

\section{讨论}
\label{sec:discuss}

本节讨论实验结果。

\end{document}
\end{lstlisting}

编译\textbf{两遍}（为什么是两遍，等下揭晓），正文里显示“如图 1 所示”“详见第 1 节”。

\textbf{逐行拆解}：

\begin{itemize}
\setlength{\itemsep}{2pt}
\item \texttt{\textbackslash{}label\{fig:result\}}：给上面那个东西（图）起个\textbf{代号}叫 \texttt{fig:result}。代号是你随便起的，用英文，别用中文；
\item \texttt{\textbackslash{}ref\{fig:result\}}：在正文里“呼叫”这个代号，LaTeX 会把它翻译成对应的\textbf{编号}；
\item 图的 \texttt{\textbackslash{}label} 和 \texttt{\textbackslash{}section} 后面的 \texttt{\textbackslash{}label}：各贴各的对象，互不干扰；
\item \texttt{\textasciitilde{}}：一个“不可断行的空格”，保证“如图 1”不被拆到两行去。写 \texttt{\textbackslash{}ref} 前面加个 \texttt{\textasciitilde{}} 是好习惯。
\end{itemize}

\textbf{灵魂操作来了}：现在把这张图删掉（或者复制一张图变成两张图），重新编译两遍——正文里的编号\textbf{自己变了}。你没有搜索，没有替换，没有顺延。这就是“编号是活的”。

代号命名的小讲究（不强制，但团队合写时救命）：

\begin{table}[htbp]
\centering
\small
\begin{tabularx}{\linewidth}{@{}*{3}{>{\raggedright\arraybackslash}X}@{}}
\toprule
\textbf{对象} & \textbf{惯用前缀} & \textbf{例子} \\
\midrule
图 & \texttt{fig:} & \texttt{fig:result} \\
表 & \texttt{tab:data} & \texttt{tab:data} \\
章节 & \texttt{sec:} & \texttt{sec:method} \\
公式 & \texttt{eq:} & \texttt{eq:loss} \\
\bottomrule
\end{tabularx}
\end{table}

前缀不是语法要求，纯是给人看的——三个月后你打开文档，\texttt{fig:cat} 和 \texttt{sec:cat} 谁是图谁是节，一眼分明。

\subsection*{5.2 编号差一的经典坑：\textbackslash{}label 必须写在 \textbackslash{}caption 之后}

这是本章最重要的坑，\textbf{按原话背下来}：\textbf{label 写在 caption 后面，写在后面，写在后面。}

看反例：

\begin{lstlisting}
\begin{figure}[htbp]
  \centering
  \includegraphics[width=0.8\textwidth]{result.png}
  \label{fig:result}   % 错！写在 caption 前面了
  \caption{实验结果对比图}
\end{figure}
\end{lstlisting}

这样编译\textbf{不报错}，正文里 \texttt{\textbackslash{}ref\{fig:result\}} 会显示“图 0”或一个错位的编号——因为 \texttt{\textbackslash{}label} 记录的是“它\textbf{上方}最近那个编号”，写在 caption 前面，它就没编号可记，记到上一个去了。

\textbf{为什么是这个顺序}：编号是 \texttt{\textbackslash{}caption} 发的，\texttt{\textbackslash{}label} 只是抄写员——抄“上面刚发出来的那个号”。caption 还没出现，它抄到的自然是别处的号。所以口诀是：\textbf{caption 发号，label 抄号；不写 caption，先发号；不发号，先写 label，必乱}。

表格同理（表在下一章细讲，这里先看结构）：

\begin{lstlisting}
\begin{table}[htbp]
  \centering
  \caption{各组实验数据}
  \label{tab:data}     % 对！在 caption 之后
  \begin{tabular}{cc}
    组别 & 准确率 \\
    A & 92\% \\
    B & 88\% \\
  \end{tabular}
\end{table}
\end{lstlisting}

\subsection*{5.3 公式引用：\textbackslash{}eqref}

公式编号的引用有个专用命令 \texttt{\textbackslash{}eqref}（在 \texttt{amsmath} 宏包里，写公式必装，下一章正式学）：

\begin{lstlisting}
\documentclass[UTF8, 12pt]{ctexart}
\usepackage{amsmath}
\begin{document}

损失函数定义如式~\eqref{eq:loss}：

\begin{equation}
  L = \frac{1}{2}(y - \hat{y})^2
  \label{eq:loss}
\end{equation}

把式~\eqref{eq:loss} 代入上式，可得……

\end{document}
\end{lstlisting}

编译两遍，正文显示“如式（1）”。\texttt{\textbackslash{}eqref} 和 \texttt{\textbackslash{}ref} 的区别就一点：\textbf{它会自动补上括号}。\texttt{\textbackslash{}ref\{eq:loss\}} 出来是 “1”，\texttt{\textbackslash{}eqref\{eq:loss\}} 出来是 “(1)”。引用公式一律用 \texttt{\textbackslash{}eqref}，省得自己打括号还打出全角半角混排。

\subsection*{5.4 “??”急救：再编译一遍}

第一次编译时，正文里的引用多半显示两个问号：\textbf{“图 ??”}。

别慌，这不是坏了，是\textbf{账还没对上}——\texttt{\textbackslash{}ref} 需要查阅 \texttt{\textbackslash{}label} 记在小账本（\texttt{.aux} 文件）里的编号，第一次编译时账本还是旧的。\textbf{再点一次排版按钮}，“??”就变成正确编号了。

这是全书第 N 次重复那条金句了，因为它真的会陪你整个 LaTeX 生涯：

\begin{notebox}
\textbf{第一次编译搭架子，第二次编译填内容。看到 “??”，再编译一遍。}
\end{notebox}

控制台里对应的警告是第 1 章见过的 \texttt{Rerun to get cross-references right}——照做即可。如果编译了两遍甚至三遍还是 “??”，那才是真出事了，多半是 \texttt{\textbackslash{}label} 和 \texttt{\textbackslash{}ref} 的代号\textbf{拼写不一致}（比如 label 写的 \texttt{fig:result}，ref 写的 \texttt{fig:results}），肉眼对着查就行。

\subsection*{5.5 复盘：写死的数字 vs 活的编号}

把两种做法摆在一起，看看代价差多少：

\begin{table}[htbp]
\centering
\small
\begin{tabularx}{\linewidth}{@{}*{3}{>{\raggedright\arraybackslash}X}@{}}
\toprule
\textbf{操作} & \textbf{Word 手打编号} & \textbf{LaTeX 自动引用} \\
\midrule
正文写"如图 3 所示" & 自己看着图打上"3" & \texttt{\textbackslash{}ref\{fig:result\}}，打代号 \\
删一张图 & 全文搜索替换十几处，还怕漏 & 什么都不做 \\
中间插一张图 & 后面全部顺延 & 什么都不做 \\
挪动图的位置 & 页码变了，编号重查 & 什么都不做 \\
被退回改稿 & 每轮都重新核一遍编号 & 编译两遍，结束 \\
\bottomrule
\end{tabularx}
\end{table}

你多付出的是什么？写 label 时多敲十几个字符。你省下的是什么？每次改稿后一两个小时的搜索替换，和一处漏改被退回的风险。这笔账，不用我帮你算了吧。

\painhead{Word 做法 vs LaTeX 做法 · 对照小结}

\begin{table}[htbp]
\centering
\small
\begin{tabularx}{\linewidth}{@{}*{3}{>{\raggedright\arraybackslash}X}@{}}
\toprule
\textbf{事项} & \textbf{Word 做法} & \textbf{LaTeX 做法} \\
\midrule
正文提到某图 & 手打"图 3" & \texttt{\textbackslash{}ref} 引用代号 \\
删图/加图 & 全文手动顺延编号 & 编译两遍，自动重排 \\
公式引用 & 手打"式（3-5）" & \texttt{\textbackslash{}eqref} 自动带括号 \\
引用失效的信号 & 没有信号，错了也不知道 & "??" 赫然在目，逼你修 \\
小账本 & 无 & 编译两遍自动对账 \\
\bottomrule
\end{tabularx}
\end{table}

\painhead{翻车合集}

\begin{table}[htbp]
\centering
\small
\begin{tabularx}{\linewidth}{@{}*{3}{>{\raggedright\arraybackslash}X}@{}}
\toprule
\textbf{症状} & \textbf{原因} & \textbf{解法} \\
\midrule
正文全是"图 ??" & 只编译了一遍 & \textbf{再编译一遍}（最后一次说这句） \\
编了两遍还是 "??" & label 和 ref 代号拼写不一致 & 全文搜代号，逐个核对大小写 \\
引用显示"图 0"或编号差一 & \texttt{\textbackslash{}label} 写在了 \texttt{\textbackslash{}caption} 前面 & label 挪到 caption 下一行 \\
\texttt{! LaTeX Error: Reference ... undefined} & \texttt{\textbackslash{}ref} 的代号在全文找不到对应 \texttt{\textbackslash{}label} & 补 label 或改对拼写 \\
引用出来没有括号 & 公式引用用了 \texttt{\textbackslash{}ref} & 换 \texttt{\textbackslash{}eqref}（记得 \texttt{\textbackslash{}usepackage\{amsmath\}}） \\
改了正文引用没更新 & 只编一遍就交了 & 交稿前固定动作：连编两遍 \\
\bottomrule
\end{tabularx}
\end{table}

\section*{练习}

\begin{enumerate}
\setlength{\itemsep}{2pt}
\item 把第 4 章的插图文档升级：给图配 \texttt{\textbackslash{}label}，正文用 \texttt{\textasciitilde{}\textbackslash{}ref\{\}} 引用它，编译两遍确认显示“图 1”；
\item 在两个 \texttt{\textbackslash{}section} 之间再插一节，给新节打 \texttt{\textbackslash{}label}，正文引用它，确认编号正确；
\item 用 \texttt{amsmath} 的 \texttt{equation} 环境写一个公式，打 \texttt{\textbackslash{}label}，正文用 \texttt{\textbackslash{}eqref} 引用；
\item 把上一题的公式\textbf{再复制一个}放在前面，编译两遍：观察两个公式和正文引用的编号怎么自动变化的；
\item \textbf{故意翻车题}：把 \texttt{\textbackslash{}label\{eq:loss\}} 挪到 \texttt{\textbackslash{}caption} 的位置思路套用到图上——故意写在 \texttt{\textbackslash{}includegraphics} 上面，编译两遍，看看 \texttt{\textbackslash{}ref} 显示出什么鬼编号。记住这个画面，以后见到它一秒确诊。
\end{enumerate}

从此，图、表、章节、公式的编号都成了活的，你的文档进入了“改稿零焦虑”状态。下一章啃 LaTeX 最大的招牌，也是 Word 最深的痛：\textbf{公式}。
